%=== CHAPTER ONE ===
%% Rewritten 2026-10-30 for the 14-state mass-only interface; 2026-11-05: the
%% deployed estimator is the reduced velocity--mass MHE (contribution 2).
%% 2026-11-06: contribution 3 leads with the deployed force-balance release rule.
%% 2026-11-06 (review 2): contributions reordered (lifecycle first); the reduced-mass
%% inertia result is stated as the modelling basis, not as a contribution.

\chapter{Introduction}
\label{ch:introduction}
\begin{spacing}{1.5}
\setlength{\parskip}{0.22in}

\section{Background}

Quadrotors are increasingly used for inspection, transport and delivery because
they can hover and manoeuvre in confined three-dimensional environments. Their
agility, however, is obtained from a light airframe with limited thrust and
torque margins. A package pickup can therefore change the vehicle dynamics on the
same time scale as the control loop. The total mass changes the
thrust-to-acceleration map, and a load hanging below the airframe multiplies the
rotational inertia that the attitude loop must move.

Model predictive control is attractive for this problem because it predicts
future motion while handling actuator limits and coupled dynamics explicitly
\cite{rawlings2017mpc,mayne2014mpc}. For a quadrotor the relevant equations are
nonlinear and the attitude evolves on the unit-quaternion manifold; nonlinear MPC
therefore avoids the loss of fidelity caused by repeated linearisation over
aggressive motion \cite{grune2017nmpc,kamel2017linear}. Its performance
nevertheless depends on the model supplied to the prediction. If the controller
continues to use the empty-airframe parameters after grasping a parcel, even a
numerically converged optimisation solves the wrong problem.

This dissertation couples moving horizon estimation with nonlinear model
predictive control. MHE estimates states and parameters by optimising over a
recent window of measurements while enforcing the system dynamics and physical
bounds \cite{rao2003}. Within that window the estimator identifies a single
payload-dependent scalar, the composite mass, and from it the controller
reconstructs every payload-dependent term in its own model.

\section{Problem Statement}

Consider a quadrotor that follows a three-dimensional figure-eight path, grasps
a payload from above, transports it and releases it at a selected phase of the
trajectory. The controller has no access to the true payload parameters. Only
ordinary onboard flight signals are available: odometry, attitude, body rate and
rotor-speed telemetry.

Two model errors appear at attachment:
\begin{enumerate}
\item the composite mass changes abruptly, shifting the hover thrust and the
translational acceleration gain;
\item the payload lever arm increases roll and pitch inertia several-fold,
reducing effective inner-loop bandwidth.
\end{enumerate}
At release both changes occur in the opposite direction. The payload is grasped
on the thrust axis, so the composite centre of mass moves only vertically and no
standing trim moment arises; a residual offset left by imperfect alignment is
treated as an unmodelled disturbance (Chapter~\ref{ch:model}).

The central research question is:

\begin{quote}
Can a quadrotor identify its payload online from its ordinary flight signals
alone --- through grasp, loaded manoeuvre and release --- and feed that estimate
to a constrained predictive controller consistently enough to remain a
well-tracking controller throughout?
\end{quote}

Release adds one decision that the continuous model does not make by itself:
when the controller may accept that a delivery has happened. Accepting a release
declares the task complete and resets auxiliary estimates learned against the
loaded vehicle; it does not change the model's mass or inertia, which follow the
estimate continuously. The controller in this work issues the release command
itself, but a command is not a detachment --- a gripper can be slow to open or
can jam. The dissertation therefore studies \emph{estimator-confirmed
detachment}: once a release has been commanded, the release is accepted only when
the estimate shows that the payload has actually gone. This costs a confirmation
delay, and it is evaluated against accepting at the command, under nominal
release and under a release whose physical detachment is held back. It is not a
detector of load loss in general: a load lost without any command is outside what
the evaluated mechanism can establish.

This question has a specific technical obstacle, which
Chapter~\ref{ch:model} addresses and which shapes the whole design. The inertia
increment is several times larger, relative to its bare-airframe value, than the
mass increment, so the rotational channel carries a strong payload signature ---
but it is also the channel this platform models least well. An estimator that
lets the mass enter its rotational model will use the mass to absorb rotational
model error. The design adopted here identifies the composite mass on the
translational channel only, derives the inertia increment from it by rigid-body
algebra, and accepts --- and measures --- the loss of information that this
costs in manoeuvre.

\section{Research Objectives}

The dissertation has five objectives:
\begin{enumerate}
\item derive a consistent quadrotor model with explicit payload mass and inertia
entry points, and reduce a centred rigid payload to one online degree of freedom
from which every payload-dependent term follows (Chapter~\ref{ch:model});
\item formulate a moving horizon estimator that identifies the composite mass
from a wrench reconstructed from measured rotor speeds, on the translational
channel, while keeping the mass from absorbing rotational model error
(Chapters~\ref{ch:model} and~\ref{ch:framework});
\item formulate a predictive controller whose prediction model, input baseline
and inner-loop command scaling are updated continuously from one health-gated
estimate (Chapter~\ref{ch:framework});
\item define a release lifecycle that separates the release command from the
physical detachment by gating release acceptance on measured evidence, and that
holds an unconfirmed release as \textsc{unresolved} rather than assuming it
complete (Chapter~\ref{ch:framework});
\item evaluate the complete interface in software-in-the-loop simulation with a
randomised, fully reported cohort and targeted ablations, and characterise where
it fails (Chapter~\ref{ch:evaluation}).
\end{enumerate}

\section{Scope}

The work covers simulation and software-in-the-loop validation of a rigidly
attached parcel grasped on the thrust axis. Airframe inertia, rotor geometry and
the thrust and
yaw-reaction coefficients are treated as calibrated plant quantities. The
calibrated airframe mass appears in the rigid-body model, in the derived payload
mass $\mP=(\mT-\mB)^{+}$ and in offline scoring.

The information boundary is stated precisely, because a looser statement would be
false. \emph{No task-specific payload mass enters the estimator or the
controller, and no attach or release notification enters the estimator.} The
controller uses its own release command only to start the confirmation timeout,
to end the attach-time inertia floor, to decide whether an acceptance also resets
the residual specific force, and --- under the command-armed and force-balance
policies --- to arm the confirmation count. The command never sets the model's mass or inertia. One mass-valued platform
constant does remain in the closed loop --- the airframe's \SI{0.30}{\kg} payload
envelope, used for an attach-time inertia floor --- and it is disclosed as a
specification of the airframe rather than as
information about the parcel. Chapter~\ref{ch:model} states exactly where it
appears.

The scope does not include eccentric grasps beyond the gripper's alignment
tolerance, deformable or sloshing loads, cable-suspended payload dynamics, wind,
motor failure, hardware flight, or a formal closed-loop stability proof for the estimator--controller
interconnection. All system-level evidence is from PX4/Gazebo
software-in-the-loop simulation. These limits are revisited in
Chapter~\ref{ch:conclusion}.

\section{Contributions}

The contributions lie in the architecture and its evaluation rather than in the
rigid-body algebra on which they rest. That basis is stated first because the
rest depends on it: for a payload grasped on the thrust axis, the parallel-axis
theorem taken about the moving composite centre of mass gives a roll and pitch
inertia increment equal to the reduced mass times the square of a nominal
vertical arm, so the composite mass is the only payload quantity to be identified
online, and the inertia follows from it and vanishes with it, with no event
needed to remove it (Chapter~\ref{ch:model}). This is a standard result applied
to the task, not a contribution in its own right. The contributions are:
\begin{enumerate}
\item \textbf{A continuous estimate-driven controller interface with an
evidence-confirmed release lifecycle.} One health-gated frame updates the
controller's prediction model, input baseline and inner-loop scaling
continuously; release acceptance is a separate decision, taken only on evidence
that the payload has left. In the deployed form that evidence is a force balance
of the thrust reconstructed from rotor speeds against the specific force measured
by the inertial unit, which does not pass through the mass estimate, and it is
opened by the release command or by a watchdog armed on airborne load. In a
$180$-flight validation of the deployed system it produced no false acceptance
during carry, never accepted a release early under a nominal thrust map,
detected every load lost without a command and accepted no jammed release. Its
price is the thrust map: when the map under-reads by more than the payload's
share of the composite mass, a loaded vehicle reads as empty, and a delayed
release was then accepted before the payload left in every such flight. In the
historical cohort, the same lifecycle on the estimator's own mass
evidence was sound only when armed by the release command: run throughout carry,
it accepted releases that had not happened during high-speed flight
(Chapter~\ref{ch:framework} and Sections~\ref{sec:eval:deployed},
\ref{sec:eval:results} and~\ref{sec:eval:campaign}).

\item \textbf{A mass estimator driven by a reconstructed physical input that
does not integrate the rotational dynamics.} The estimator is fed the thrust
reconstructed from measured rotor speeds rather than the controller's intended
input, which would make the mass unobservable by construction. In its deployed
form it keeps only the world-frame velocity and the composite mass as states and
takes the measured attitude as a known input, so that errors of the poorly
modelled rotational channel have no path to the mass; a mass-health check flags
an estimate that leaves a physical band instead of hiding it behind a hard bound.
Replayed on the 312 recorded campaign flights, it had a smaller loaded mass error
than the full-state estimator in every nominal flight. The evaluation campaigns
were flown with that earlier full-state form, which separated the channels more
weakly by carrying the payload inertia as a parameter lagged by one solve. In a
paired closed-loop comparison of $96$ flights with everything else fixed, the
deployed form met every pre-specified acceptance criterion: it lowered the
loaded mass error from \SI{0.035}{\kg} to \SI{0.013}{\kg}, removed the
manoeuvre-induced collapse of the estimate, solved about ten times faster and
changed neither tracking nor the lifecycle. Replay against an extended Kalman
filter and simpler estimators on the same inputs shows that this gain comes from
the model structure --- reconstructed thrust and measured attitude as known
inputs, no rotational channel --- rather than from the moving-horizon solution,
which matched a filter on accuracy (Chapters~\ref{ch:model},
\ref{ch:framework} and Section~\ref{sec:eval:reduced}).

\item \textbf{A software-in-the-loop evaluation with randomised blocks, complete
reporting, ablations and failure cases.} A $180$-flight validation of the
deployed system from a tagged commit, with acceptance criteria inherited from an
earlier pre-specification; and, as the historical record of how the design was
reached, a frozen cohort of $120$ flights and a release campaign of $312$, both
with the full-state estimator, and a paired estimator comparison of $96$. Every
flight is reported, with the configuration each
result applies to stated explicitly, together with an
ablation of the residual specific force and an account of the failure modes --- in particular the
collapse of the mass estimate under sustained manoeuvre and the false
acceptances it causes in an always-on lifecycle, and, for the deployed system,
the early acceptance caused by an under-reading thrust map
(Chapter~\ref{ch:evaluation}).
\end{enumerate}

\section{Dissertation Organisation}

Chapter~\ref{ch:literature} reviews predictive control, online estimation,
inertial-parameter identification and payload-adaptive aerial control, and
positions this work against representative approaches.
Chapter~\ref{ch:model} derives the model and the payload parameter structure, and
is where the one-degree-of-freedom reduction and the channel assignment of the
mass are developed in full.
Chapter~\ref{ch:framework} presents the coupled architecture, the numerical form
of both optimal control problems, the continuous payload interface and the
release lifecycle. Chapter~\ref{ch:evaluation} describes the implementation and
reports the validation of the deployed system and the historical experiments
that led to it, including the conditions under which
the interface fails. Chapter~\ref{ch:conclusion} answers the research objectives
against that evidence and states what remains open.

\end{spacing}
\newpage
